\documentclass[a4paper,11pt]{article}
\usepackage[T1]{fontenc}
\usepackage{lmodern}
\usepackage{amsmath,amssymb,amsthm}
\usepackage{booktabs}
\usepackage[a4paper,hmargin=32mm,top=26mm,bottom=27mm]{geometry}
\usepackage[unicode,hidelinks]{hyperref}
\hypersetup{
  pdftitle={Equality in a binomial gcd bound of Erdős and Szekeres},
  pdfauthor={Yoshiharu Onishi},
  pdfsubject={Binomial coefficients, prime-power divisors and equality in Erdős Problem 700(a)},
  pdfkeywords={binomial coefficients, Kummer's theorem, digit sums, Lean},
  pdflang={en}
}

\newtheorem{theorem}{Theorem}[section]
\newtheorem{lemma}[theorem]{Lemma}
\newtheorem{proposition}[theorem]{Proposition}
\newtheorem{corollary}[theorem]{Corollary}
\newcommand{\colg}[2]{g_{#1}(#2)}
\newcommand{\crit}{\mathcal D}
\newcommand{\ds}{\operatorname{s}}
\newcommand{\ord}{\operatorname{ord}}
\newcommand{\lcm}{\operatorname{lcm}}
\newcommand{\leanname}[1]{{\small\texttt{\detokenize{#1}}}}
\setcounter{tocdepth}{2}
\setlength{\emergencystretch}{1em}
\setlength{\belowcaptionskip}{5pt}
\renewcommand{\arraystretch}{1.12}
\linespread{1.045}
\clubpenalty=8500
\widowpenalty=8500
\raggedbottom

\title{\bfseries Equality in a binomial gcd bound\\of Erd\H{o}s and Szekeres}
\author{Yoshiharu Onishi}
\date{November 20, 2026}

\begin{document}
\maketitle

\begin{abstract}
For a composite integer \(n\), let \(f(n)\) be the minimum of
\(\gcd(n,\binom nk)\) over \(2\le k\le n/2\). Erd\H{o}s and Szekeres asked
which \(n\) satisfy \(f(n)=n/Q(n)\), where \(Q(n)\) is the largest prime
power dividing \(n\). We give a necessary and sufficient condition using
intersections of subset sums obtained by splitting the prime-power terms in
the base expansions of \(n\). It also determines equality when the
denominator is the largest prime factor. For each fixed \(m\ge2\), equality
\(f(mp)=m\) has relative density one among primes \(p>m\). For every prime
\(q\nmid m\), the exponents \(e\) with \(f(mq^e)=m\) have density one.
Both exceptional sets satisfy power-saving bounds. We give exact formulas
for several families, including one in which the exceptions correspond to
powers of two that omit the ternary digit \(2\). The main theorems have
been formalized in Lean~4.
\end{abstract}

\section{The two equality questions}\label{sec:intro}

For a composite integer \(n\ge4\), write
\begin{equation}\label{eq:definition}
  \colg nk=\gcd\!\left(n,\binom nk\right),\qquad
  f(n)=\min_{2\le k\le\lfloor n/2\rfloor}\colg nk .
\end{equation}
The largest prime factor \(P(n)\) and the largest prime-power divisor
\(Q(n)\) are given by
\[
  P(n)=\max_{p\mid n}p,\qquad Q(n)=\max_{p\mid n}p^{v_p(n)}.
\]
In both maxima \(p\) runs over primes, and \(v_p\) is the \(p\)-adic
valuation. Erd\H{o}s and Szekeres asked for the composite integers with
\(f(n)=n/Q(n)\) \cite[pp.~98--99]{ES}. Problem~700(a) in Bloom's collection
uses \(P(n)\) \cite{Bloom}. For \(n=12\) the five gcds, in increasing order
of \(k\), are \(6,4,3,12,12\). Thus \(f(12)=3=12/Q(12)\), whereas
\(12/P(12)=4\). The discussion accompanying Problem~700 also describes
digit conditions for the family \(n=q(q+1)\) with \(q\) prime \cite{Bloom}.

Theorem~\ref{thm:criterion} gives a finite test for both equalities from
the prime factorization and the digits of \(n\). The argument also gives a
column witnessing strict inequality whenever an intersection contains an
interior point. Sections~\ref{sec:primes} and~\ref{sec:powers} use carries in
base expansions to count exceptions with fixed cofactor. Explicit
families in Section~\ref{sec:families} show how the exceptional set can
be described exactly. For instance, when \(p>8\) is prime, \(f(8p)=4\)
precisely when \(p-1\) or \(3p-1\) is a power of two. The present
characterization concerns part~(a). Parts~(b) and~(c) ask, respectively,
whether \(f(n)>\sqrt n\) occurs infinitely often and how large \(f(n)\)
can be relative to \(n/(\log n)^\alpha\).

\section{A criterion for every composite integer}\label{sec:criterion}

Fix \(n\), and abbreviate \(P=P(n)\) and \(Q=Q(n)\). For \(d>1\), let
\(s(d)\) be its smallest prime factor. Use the divisors
\begin{equation}\label{eq:critical}
  \crit(n)=\{d\mid n:\ d>Q,\ d/s(d)\le Q\}.
\end{equation}
Given a prime \(p\mid n\), write \(n=\sum_j n_jp^j\), with
\(0\le n_j<p\). Start with \(n_j\) copies of each \(p^j\). One
\emph{split} replaces a single \(p^j\), \(j\ge1\), by \(p\) copies of
\(p^{j-1}\). For \(1\le b\le v_p(n)\), let \(S_{p,b}(n)\) be the set of
all subset sums of the resulting multiset, allowing any sequence of at most
\(v_p(n)-b\) splits. The empty and full sums give \(0,n\in S_{p,b}(n)\).
For \(d>1\) dividing \(n\), write
\[
  I_d(n)=\bigcap_{p^b\parallel d}S_{p,b}(n),
\]
where \(p^b\parallel d\) ranges over the exact prime powers of \(d\), with
\(b\ge1\).

\begin{theorem}\label{thm:criterion}
Suppose \(n\) has at least two distinct prime factors. Then
\begin{align}
 f(n)=n/Q &\ \Longleftrightarrow\
      I_d(n)=\{0,n\}\quad\text{for every }d\in\crit(n),\label{eq:Q}\\
 f(n)=n/P &\ \Longleftrightarrow\
      P=Q\ \text{ and }\ I_d(n)=\{0,n\}\quad\text{for every }d\in\crit(n).
      \label{eq:P}
\end{align}
For \(n=p^a\) with \(a\ge2\), one has \(f(n)=p\). Among prime powers the
first equality has no solutions, and the second holds precisely for \(a=2\).
\end{theorem}

\begin{lemma}\label{lem:splits}
For \(0\le k\le n\), the least number of splits required to realize \(k\)
as a subset sum in base \(p\) is \(v_p\binom nk\).
\end{lemma}

\begin{proof}
Write \(\ds_p(x)\) for the sum of the base-\(p\) digits of \(x\). The
canonical expansion uses \(\ds_p(n)\) terms. Each split adds \(p-1\) terms.
If a submultiset has sum \(k\), it contains at least \(\ds_p(k)\) terms,
since combining \(p\) equal powers into the next power reduces the count
until the canonical expansion is reached. Its complement similarly needs
at least \(\ds_p(n-k)\) terms. Conversely, combine the canonical expansions
of \(k\) and \(n-k\) and perform such carries until their union is the
canonical expansion of \(n\). Reversing those carries attains the lower
bound. The least number of splits is therefore
\begin{equation}\label{eq:kummer}
 \frac{\ds_p(k)+\ds_p(n-k)-\ds_p(n)}{p-1}=v_p\binom nk .
\end{equation}
The equality on the right is Kummer's digit-sum formula \cite{Granville}.
\end{proof}

\begin{proof}[Proof of Theorem~\ref{thm:criterion}]
First suppose \(n=mp^a\), where \(p\nmid m\) and \(m\ge2\).
The identity
\(\binom{n}{p^a}=m\binom{n-1}{p^a-1}\) gives divisibility by \(m\).
Comparing coefficients of \(X^{p^a}\) in
\((1+X)^n\equiv(1+X^{p^a})^m\pmod p\) shows that this coefficient is
\(m\bmod p\). Consequently
\begin{equation}\label{eq:attain}
 \colg n{p^a}=m.
\end{equation}
Taking \(p^a=Q\) gives an admissible column attaining \(n/Q\).

For any \(0\le k\le n\), let \(g=\colg nk\). Then \(g<n/Q\) is equivalent
to \(n/g>Q\). In that case choose the least divisor \(d\) of \(n/g\) that
exceeds \(Q\). Its largest proper divisor is \(d/s(d)\), so \(d\in\crit(n)\).
Conversely, any \(d\in\crit(n)\) dividing \(n/g\) gives \(g<n/Q\).
At a prime power \(p^b\parallel d\), put \(a=v_p(n)\). The condition
\(p^b\mid n/g\) is precisely \(v_p\binom nk\le a-b\). By
Lemma~\ref{lem:splits} this is equivalent to \(k\in S_{p,b}(n)\).
Thus \(d\mid n/g\) is equivalent to \(k\in I_d(n)\).

The endpoints belong to every \(I_d(n)\). At \(k=1,n-1\) the gcd equals
\(n\), and symmetry interchanges \(k\) and \(n-k\). Hence every interior
point of an intersection for \(d\in\crit(n)\) yields a witness in the
range of \eqref{eq:definition}. This proves \eqref{eq:Q}. Since
\(P\le Q\) and \(f(n)\le n/Q\), equality with \(n/P\) forces \(P=Q\).
This proves \eqref{eq:P}.

Finally, for \(n=p^a\) and \(0<k<p^a\) one has
\(v_p\binom{p^a}{k}=a-v_p(k)\). Indeed
\(\binom{p^a}{k}=(p^a/k)\binom{p^a-1}{k-1}\), and the second binomial is
a unit modulo \(p\) by Lucas's theorem \cite{Granville}.
The valuation is at least one and is one at \(k=p^{a-1}\).
Now \(n/Q=1\) and \(n/P=p^{a-1}\), proving the last assertion.
\end{proof}

\paragraph{Examples.}
For \(n=30\), one has \(Q=5\) and \(\crit(30)=\{6,10,15\}\).
No splits are available. The base-two terms are \(2,4,8,16\), so their
subset sums are even. The interior base-three sums are \(3,27\), and the
interior base-five sums are \(5,25\). These three interior sets are
pairwise disjoint, giving \(f(30)=6\), an example in \cite{ES}.
For \(n=45\), one has \(Q=9\) and \(\crit(45)=\{15\}\).
In base three, split one \(9\) in \(27+9+9\) to obtain the subsum
\(9+3+3=15\). In base five, three of the four \(5\)'s in
\(25+5+5+5+5\) give the same sum. Hence \(15\mid45/\colg{45}{15}\),
so this column is at most \(3\). The identity
\(n\mid k\binom nk\) implies \(n\mid k\colg nk\), by cancellation after
dividing both \(n\) and \(\binom nk\) by their gcd. In particular no
interior gcd is \(1\). The least possible divisor of \(45\) greater than
\(1\) is \(3\), giving \(f(45)=3<5\).

For squarefree \(n\) no splits are possible. Membership in \(S_{p,1}(n)\)
then means that each base-\(p\) digit of \(k\) is at most the corresponding
digit of \(n\). For \(n=10815=3\cdot5\cdot7\cdot103\), the divisor
\(105\) is in \(\crit(n)\). The digits of \(k=3255\) satisfy this test
in all three bases \(3,5,7\), giving \(f(n)\le n/105=103<105=n/Q\).

\paragraph{Computing the intersections.}
The sets \(S_{p,b}(n)\) can be computed without listing refinements.
Read the base-\(p\) digits of a candidate \(k\) from least to most
significant, subtracting them from \(n\). With incoming borrow \(u\in\{0,1\}\),
digit \(x\) at position \(i\) produces a new borrow precisely when
\(x+u>n_i\). Keep the partial sums \(\sum_{j\le i}x_jp^j\) for each
pair consisting of the outgoing borrow and the number of borrows so far.
Start at sum zero with no borrow, discard counts above \(v_p(n)-b\),
and require zero outgoing borrow at the final digit. Subtraction borrows
are the carries in \(k+(n-k)\), so \eqref{eq:kummer} proves that the
surviving sums are precisely \(S_{p,b}(n)\). Bitsets indexed by \(k\)
then implement the intersections in \eqref{eq:Q}.

\section{A fixed cofactor and a varying prime}\label{sec:primes}

When \(n=mp\) with \(p>m\) prime, one has \(P(n)=Q(n)=p\).
The problem is to decide whether some column improves on the value \(m\)
at \(k=p\). We first record a divisibility test that works simultaneously
at all columns of a scaled row.

\begin{lemma}\label{lem:block}
Let \(\ell\) be prime, \(Y\ge A\ge2\), and \(a\ge1\). Put
\(H=\lceil\log_\ell Y\rceil\) and \(L=H+a-1\). Suppose that the
base-\(\ell\) expansion of a positive integer \(z\) contains \(L\)
consecutive digits equal to \(\ell-1\). Then
\[
   \ell^a\mid\binom{Az}{tz}\qquad(1\le t<A).
\]
\end{lemma}

\begin{proof}
Let the run start at position \(d\), with the units digit at position
zero, and put \(v=\ell^d-(z\bmod\ell^d)\). Thus \(1\le v\le\ell^d\).
For \(H\le j\le H+a-1\), put \(M=\ell^{d+j}\). The run gives
\(z\equiv-v\pmod M\) and \(Av\le M\). The least nonnegative residues
of \(tz\) and \((A-t)z\) are \(M-tv\) and \(M-(A-t)v\).
Their sum \(2M-Av\) is at least \(M\), so addition produces a carry
into position \(d+j\). The \(a\) choices of \(j\) give distinct carries.
Kummer's theorem gives the required divisibility.
\end{proof}

\begin{theorem}\label{thm:prime}
Fix \(m\ge2\), and let \(E_m(X)\) count the primes \(m<p\le X\) for
which \(f(mp)<m\). There are constants \(C_m,\delta_m>0\) such that
\(E_m(X)\le C_m X^{1-\delta_m}\) for every integer \(X\ge1\).
Consequently \(f(mp)=m\) on a set of primes of relative density one.
\end{theorem}

\begin{proof}
Since \(mp\mid k\colg{mp}{k}\), each column with \(p\nmid k\) has gcd
at least \(p>m\). At \(k=tp\), \(1\le t\le m/2\), Lucas's theorem gives
\(\binom{mp}{tp}\equiv\binom mt\not\equiv0\pmod p\). Thus the gcd at
this column is \(\gcd(m,\binom{mp}{tp})\), and \eqref{eq:attain} gives
equality at \(t=1\). We have proved
\begin{equation}\label{eq:primecolumns}
 f(mp)=m\quad\Longleftrightarrow\quad
 m\mid\binom{mp}{tp}\quad(1\le t\le\lfloor m/2\rfloor).
\end{equation}
For each \(\ell^a\parallel m\), take \(Y=A=m\) in Lemma~\ref{lem:block}
and write \(B_\ell=\ell^{\lceil\log_\ell m\rceil+a-1}\).
One occurrence of the digit \(B_\ell-1\) in base \(B_\ell\) supplies a
whole run of maximal base-\(\ell\) digits. If this happens for every
\(\ell\mid m\), all the divisibilities in \eqref{eq:primecolumns} follow.

For \(B=B_\ell\), the integers \(0\le z\le X\) can be written with
\(r=1+\lfloor\log_B X\rfloor\) base-\(B\) digits, padding with zeros.
At most \((B-1)^r\) such strings omit \(B-1\). Summing over the primes
dividing \(m\) gives
\[
 E_m(X)\le\sum_{\ell\mid m}(B_\ell-1)^{1+\lfloor\log_{B_\ell}X\rfloor}.
\]
For \(B\ge2\) set \(\theta(B)=\log(B-1)/\log B<1\). Each summand
is at most \((B_\ell-1)X^{\theta(B_\ell)}\). This proves the bound
with \(\delta_m=1-\max_{\ell\mid m}\theta(B_\ell)>0\).
Chebyshev's lower bound \(\pi(X)\gg X/\log X\), where \(\pi\) counts
primes, gives \(E_m(X)/\pi(X)\to0\).
\end{proof}

\section{A fixed cofactor and a varying exponent}\label{sec:powers}

Fix \(m\ge2\) and a prime \(q\nmid m\). Equation~\eqref{eq:attain}
gives \(f(mq^e)\le m\) for all \(e\ge1\). For sufficiently large
\(e\), the prime power \(q^e\) is \(Q(mq^e)\), so equality with \(m\)
is the original Erd\H{o}s--Szekeres equality. We first bound the number
of candidate columns independently of \(e\).

\begin{proposition}\label{prop:candidates}
For \(e\ge1\), put
\[
 h=\min\{e,\lfloor\log_q(m-1)\rfloor\},\quad z=q^{e-h},\quad A=mq^h.
\]
Then \(A\le m(m-1)\). Every column with gcd below \(m\) in the range
of \eqref{eq:definition} has the form \(k=tz\), \(1\le t\le A/2\).
The gcd at each such index is
\begin{equation}\label{eq:powercolumns}
 \colg{Az}{tz}
 =q^{\min\{e,v_q\binom At\}}
   \prod_{\ell^a\parallel m}\ell^{\min\{a,v_\ell\binom{Az}{tz}\}} .
\end{equation}
\end{proposition}

\begin{proof}
The bound on \(A\) follows from \(q^h\le m-1\). If \(h=e\) then
\(z=1\). Otherwise \(m\le q^{h+1}\). For \(g=\colg{Az}{k}<m\)
one has \(v_q(g)\le h\), and \(q^e\mid kg\) forces \(q^{e-h}\mid k\).
The bound on \(t\) is now the admissible half-row bound. Multiplication
by \(z\) shifts base-\(q\) digits, preserving the carries in
\(t+(A-t)\). Thus \(v_q\binom{Az}{tz}=v_q\binom At\), and factoring
the gcd yields \eqref{eq:powercolumns}. If \(tz=1\), then
\(\colg{Az}{1}=Az\), so including this index in a test for \(g<m\)
does not affect the result.
\end{proof}

\begin{theorem}\label{thm:exponents}
For every \(m\ge2\) there exists \(\eta_m>0\) such that for every prime
\(q\nmid m\) there is a constant \(C_{m,q}>0\) with
\begin{equation}\label{eq:exponents}
 \#\{1\le e\le E: f(mq^e)<m\}
 \le C_{m,q}E^{1-\eta_m}\qquad(E\ge1).
\end{equation}
In particular \(f(mq^e)=m\) on a set of exponents of natural density one.
\end{theorem}

\begin{proof}
Set \(h_0=\lfloor\log_q(m-1)\rfloor\), \(Y=m(m-1)\). For \(e\ge h_0\)
write \(s=e-h_0\). For each \(\ell^a\parallel m\) use
\[
 L_\ell=\lceil\log_\ell Y\rceil+a-1,\qquad B_\ell=\ell^{L_\ell}.
\]
If the expansion of \(q^s\) in base \(B_\ell\) contains \(B_\ell-1\)
for every \(\ell\mid m\), Lemma~\ref{lem:block} with \(A=mq^{h_0}\)
gives \(m\mid\binom{Aq^s}{tq^s}\) for all \(1\le t<A\).
Proposition~\ref{prop:candidates} then gives equality.

We bound how often an expansion can omit its maximal digit. For each
prime \(\ell\ne q\) there exists \(c=c_{\ell,q}\ge0\) such that
\begin{equation}\label{eq:periodloss}
 q^u\equiv q^v\pmod{\ell^R}
 \quad\Longrightarrow\quad
 u\equiv v\pmod{\ell^{\max(R-c,0)}} .
\end{equation}
For \(d=|u-v|>0\) the first congruence implies \(\ell^R\mid q^d-1\).
If \(\ell\) is odd, apply the valuation formula
\(v_\ell(x^d-1)=v_\ell(x-1)+v_\ell(d)\) to \(x=q^{\ell-1}\equiv1\pmod\ell\).
Since \(q^d-1\mid x^d-1\), this bounds \(R\).
If \(\ell=2\), \(q\) is odd and
\(v_2(q^{2d}-1)=v_2(q-1)+v_2(q+1)+v_2(d)\), which gives the same
conclusion. Indeed, for \(\ell\mid y\) and \(r\) prime to \(\ell\), the
binomial expansion gives \(v_\ell((1+y)^r-1)=v_\ell(y)\). For \(r=\ell\)
the valuation increases by one if \(\ell\) is odd, or if \(\ell=2\)
and \(4\mid y\). In the even case start with \(1+y=q^2\).
If \(u=v\), \eqref{eq:periodloss} is immediate.

Fix \(\ell\mid m\), set \(B=B_\ell\), \(r=1+\lfloor\log_B E\rfloor\),
and take \(R=L_\ell r\). Then \(E<B^r=\ell^R\).
The bottom \(r\) digits of \(q^s\) have at most \((B-1)^r\) possible
values if they omit \(B-1\). By \eqref{eq:periodloss}, exponents
yielding any one such residue belong to a single class modulo
\(\ell^{\max(R-c,0)}\). Since \(0\le s\le E<\ell^R\), at most
\(\ell^c\) exponents give that residue. The total exceptional count is
therefore bounded by
\[
 h_0+\sum_{\ell\mid m}\ell^{c_{\ell,q}}
      (B_\ell-1)^{1+\lfloor\log_{B_\ell} E\rfloor}.
\]
As in Theorem~\ref{thm:prime}, take
\(\eta_m=1-\max_{\ell\mid m}\theta(B_\ell)>0\). The bases \(B_\ell\)
depend only on \(m\) and its prime factors, giving the asserted uniformity
in \(q\). Dividing \eqref{eq:exponents} by \(E\) proves the density statement.
\end{proof}

\begin{corollary}\label{cor:progression}
Suppose \(e_0\ge\max\{1,\lfloor\log_q(m-1)\rfloor\}\) and \(f(mq^{e_0})=m\).
For the primes \(\ell\mid m\) put
\[
 R_\ell=1+\lfloor\log_\ell(mq^{e_0})\rfloor,\qquad
 T=\lcm_{\ell\mid m}\ord_{\ell^{R_\ell}}(q),
\]
where \(\ord_N(q)\) is the multiplicative order of \(q\) modulo \(N\).
Then \(f(mq^{e_0+jT})=m\) for all \(j\ge0\). Every residue class modulo
any \(D\ge1\) contains such an infinite progression of equality exponents,
with arbitrarily large first term.
\end{corollary}

\begin{proof}
The value \(h=\lfloor\log_q(m-1)\rfloor\) is fixed above \(e_0\).
Put \(z_0=q^{e_0-h}\) and \(z=q^{e_0+jT-h}\). Then
\(z\equiv z_0\pmod{\ell^{R_\ell}}\). For each candidate \(t\) in
Proposition~\ref{prop:candidates}, the residues of \(tz,(A-t)z\) agree
through the lowest \(R_\ell\) digits with those at \(z_0\).
Every old carry occurs within those digits because \(Az_0<\ell^{R_\ell}\).
Hence the \(\ell\)-adic valuation in \eqref{eq:powercolumns} cannot
decrease. Its \(q\)-part also cannot decrease because \(v_q\binom At\)
is fixed. Each old gcd thus divides the corresponding new gcd. All were
at least \(m\), and \eqref{eq:attain} still gives \(m\).
Finally each residue class modulo \(D\) has positive density, so
Theorem~\ref{thm:exponents} supplies an arbitrarily large seed \(e_0\)
there. Replacing \(T\) by \(\lcm(T,D)\) keeps the progression in the class.
\end{proof}

\newpage
\section{Exact values for small cofactors}\label{sec:families}

\subsection{Cofactor eight and a varying prime}

\begin{proposition}\label{prop:eight}
For a prime \(p>8\), one has \(f(8p)=4\) if \(p-1\) or \(3p-1\) is
a power of two. In all other cases \(f(8p)=8\).
\end{proposition}

\begin{proof}
Write \(w(x)=\ds_2(x)\). Only \(k=tp\) with \(1\le t\le4\) can improve
on \(8\), and the prime \(p\) drops out of each gcd as in
\eqref{eq:primecolumns}. For \(t=1,3\), the relation
\(8\mid tp\,\colg{8p}{tp}\) forces the gcd to be \(8\).
For \(t=2\), the valuation in \eqref{eq:kummer} is
\(w(2p)+w(6p)-w(8p)=w(3p)\). For \(t=4\) it is \(2w(4p)-w(8p)=w(p)\).
Consequently
\[
  f(8p)=2^{\min\{3,w(p),w(3p)\}} .
\]
An odd integer \(x>1\) has at least two binary ones. It has exactly two
precisely when \(x=2^j+1\) for some \(j\ge1\).
\end{proof}

For example, \(3\cdot11-1=32\) gives \(f(88)=4\), while \(f(104)=8\).
Each exponent \(j\) gives at most two primes from the two equations, so
this family has only \(O(\log X)\) exceptions up to \(p\le X\).

\subsection{Cofactors four and six at all exponents}

\begin{theorem}\label{thm:small}
For every odd prime \(q\) and \(e\ge1\), \(f(4q^e)=4\) apart from
\(f(12)=3\). For every \(e\ge1\), \(f(6\cdot5^e)=6\).
\end{theorem}

\begin{proof}
For \(q\ge5\), put \(x=q^e\). In the row \(4x\),
Proposition~\ref{prop:candidates} leaves \(k=x,2x\). The first gives
\(4\) by \eqref{eq:attain}. At the centre,
\(v_2\binom{4x}{2x}=2w(2x)-w(4x)=w(x)\ge2\), since \(x>1\) is odd.
Hence the second gcd is a multiple of \(4\).

For \((m,q)=(4,3),(6,5)\) and \(e\ge2\), take \(h=1\), \(A=mq\),
\(z=q^{e-1}\). If \(g=\colg{Az}{tz}<m\), then \(A\mid tg\), by
\(Az\mid tz\,g\), and \eqref{eq:powercolumns} forces
\[
 \min\{2,v_q(g)\}=\min\!\left\{2,v_q\binom At\right\}.
\]
Table~\ref{tab:exclusions} lists all \(1\le t\le A/2,\ 1\le g<m\)
with \(A\mid tg\). The valuations in its last column follow from
\eqref{eq:kummer}. Each pair fails the displayed equality.
Thus \eqref{eq:attain} gives the asserted value. For \(e=1\), the
values \(f(12)=3\) and \(f(30)=6\) were computed in
Sections~\ref{sec:intro} and~\ref{sec:criterion}.
\end{proof}

\begin{table}[ht]
\centering
\caption{Candidate pairs in the two remaining rows, listed in matching order.}
\label{tab:exclusions}
\begin{tabular}{@{}cccll@{}}
\toprule
\(m\)&\(q\)&\(A\)& \((t,g)\) with \(A\mid tg\) & \(v_q\binom At\)\\
\midrule
4&3&12& \((6,2),(4,3)\) & \(1,2\)\\
6&5&30& \((15,2),(10,3),(15,4),(6,5),(12,5)\) & \(1,1,1,2,2\)\\
\bottomrule
\end{tabular}
\end{table}

\clearpage
\subsection{Cofactor nine and ternary digits}

\begin{theorem}\label{thm:nine}
For \(e\ge4\),
\[
 f(9\cdot2^e)=
 \begin{cases}
  4,&\text{if the ternary expansion of \(2^{e-1}\) omits the digit \(2\)},\\
  9,&\text{if that expansion contains the digit \(2\)}.
 \end{cases}
\]
In particular the set of \(e\ge4\) giving the value \(4\) is finite if
and only if the set of \(t\ge0\) for which \(2^t\) omits the ternary
digit \(2\) is finite.
\end{theorem}

\begin{proof}
Set \(n=9\cdot2^e\). The column \(k=2^e\) gives \(9\) by
\eqref{eq:attain}. Suppose an admissible \(k\) has \(g=\colg nk<9\).
If \(3\nmid k\), then \(9\mid g\) by \(n\mid kg\). If \(3\mid k\)
but \(9\nmid k\), the same relation gives \(3\mid g\), so \(v_2(g)\le1\)
and \(2^{e-1}\mid k\). Writing \(k=t2^{e-1}\) leaves \(t=3,6\).
Deleting a common binary shift, the binomial valuations
\(v_2\binom{18}{3},v_2\binom{18}{6}\) are \(4,2\). Both contradict
\(v_2(g)\le1\) since \(e\ge4\).

It remains to consider \(9\mid k\). Now \(v_2(g)\le3\), so
\(k=9\cdot2^{e-3}u\) with \(1\le u\le4\). The four values
\(v_2\binom{72}{9u}\) are \(6,4,6,2\), respectively, by
\eqref{eq:kummer}. Only \(u=4\) remains, namely \(k=n/2\), and its
binomial coefficient has 2-adic valuation two. In base three, adding
\(n/2=9\cdot2^{e-1}\) to itself produces a carry precisely when its
expansion has a digit \(2\). Multiplication by \(9\) appends two zeros.
Thus at the centre the gcd is \(4\) if \(2^{e-1}\) omits \(2\),
and is divisible by \(12\) otherwise. This proves the formula. The
shift \(t=e-1\) leaves only \(t=0,1,2\) outside the comparison, which
proves the finiteness equivalence.
\end{proof}

The finiteness question in Theorem~\ref{thm:nine} is Erd\H{o}s
Problem~406 \cite{digits}. Saye \cite{Saye} verified that, for
\(0\le t\le2\cdot3^{45}\), only \(2^t=1,4,256\) omit the ternary
digit \(2\). Thus the unique exceptional exponent in the verified
range \(4\le e\le2\cdot3^{45}+1\) is \(e=9\), for which
\(2^8=(100111)_3\) and \(f(4608)=4\). The formula leaves open the
finiteness of the full exceptional set. Holdum, Klausen and Rasmussen
\cite{HKR} used the digits of powers in prime bases to bound exceptions
to divisibility for central binomial coefficients. Here the reduction
to a central column follows from the exclusion of all other possible
minimizers below \(9\).

\appendix
\section{Formal proofs and numerical checks}\label{sec:verification}

Lean~4 and mathlib are pinned to \texttt{v4.33.0-rc1}.
Table~\ref{tab:lean} lists the main declarations in the namespace
\texttt{Erdos700}. The theorem \leanname{equality_sets} uses the infimum
definition of \(f\) in the Formal Conjectures statement \cite{formal}.

\begin{table}[!ht]
\centering
\caption{Main results and their Lean declarations.}
\label{tab:lean}
\small
\renewcommand{\arraystretch}{1.02}
\begin{tabular}{@{}ll@{}}
\toprule
Result & Declaration in \texttt{Erdos700}\\
\midrule
Theorem~\ref{thm:criterion} & \leanname{all_composite_criteria}, \leanname{equality_sets}\\
Theorem~\ref{thm:prime} & \leanname{exceptional_prime_power_saving}\\
Theorem~\ref{thm:exponents} & \leanname{exceptional_exponent_power_saving}\\
Corollary~\ref{cor:progression} & \leanname{equality_progression_of_seed}\\
 & \leanname{equality_progression_in_every_class}\\
Proposition~\ref{prop:eight} & \leanname{eight_prime_classification}\\
Theorem~\ref{thm:small} & \leanname{four_odd_prime_power}, \leanname{six_five_power}\\
Theorem~\ref{thm:nine} & \leanname{nine_mul_two_pow}, \leanname{nine_exceptions_finite_iff}\\
\bottomrule
\end{tabular}
\end{table}

\paragraph{Numerical checks.}
Two Python programs accompany the proofs in the author's repository
(\href{https://github.com/yoshiyoyoyo44/erdos700}{\nolinkurl{github.com/yoshiyoyoyo44/erdos700}}).
The first compares
Theorem~\ref{thm:criterion} with the minimum computed from exact
binomial coefficients for all \(4{,}330\) composites \(n\le5{,}000\).
There are \(2{,}811\) equalities with \(n/Q(n)\) and \(2{,}578\) with
\(n/P(n)\). For \(n\le144\), it also enumerates the refinements directly.

The second program checks Proposition~\ref{prop:eight} against the
cofactor formula for all \(9{,}588\) primes \(8<p\le100{,}000\), with
direct full-row checks for the \(91\) primes \(p\le500\) in that range.
The exceptional primes found were
\[
 11,\quad17,\quad43,\quad257,\quad683,\quad2731,\quad43691,\quad65537.
\]
It also checks \(19{,}840\) instances of Theorem~\ref{thm:small},
including \(130\) direct row minima. The repository README gives the
commands for reproducing the computations and checking the formal proofs.
OpenAI Codex assisted with the English manuscript and the Lean formalization.

\pdfbookmark[1]{References}{references}
\begin{thebibliography}{9}
\bibitem{ES}
P.~Erd\H{o}s and G.~Szekeres,
\href{https://static.renyi.hu/~p_erdos/1978-46.pdf}{Some number theoretic problems on binomial coefficients}.
\emph{Aust. Math. Soc. Gaz.} \textbf{5} (1978), 97--99.

\bibitem{Bloom}
T.~F.~Bloom (ed.), \href{https://www.erdosproblems.com/700}{Erd\H{o}s Problem 700},
\emph{Erd\H{o}s Problems}. See also the
\href{https://www.erdosproblems.com/forum/thread/700?order=oldest}{discussion of special families}.

\bibitem{Granville}
A.~Granville,
\href{https://www.cecm.sfu.ca/organics/papers/granville/}{Arithmetic properties of binomial coefficients I: Binomial coefficients modulo prime powers}.
\emph{CMS Conf. Proc.} \textbf{20} (1997), 253--275.

\bibitem{digits}
T.~F.~Bloom (ed.), \href{https://www.erdosproblems.com/406}{Erd\H{o}s Problem 406},
\emph{Erd\H{o}s Problems}.

\bibitem{Saye}
R.~I.~Saye,
\href{https://cs.uwaterloo.ca/journals/JIS/VOL25/Saye/saye3.html}{On two conjectures concerning the ternary digits of powers of two}.
\emph{J. Integer Seq.} \textbf{25} (2022), Article 22.3.4.

\bibitem{HKR}
S.~T.~Holdum, F.~R.~Klausen and P.~M.~R.~Rasmussen,
\href{https://math.colgate.edu/~integers/p43/p43.pdf}{Powers in prime bases and a problem on central binomial coefficients}.
\emph{Integers} \textbf{15} (2015), A43.

\bibitem{formal}
Google DeepMind and contributors,
\href{https://github.com/google-deepmind/formal-conjectures/blob/2c24ebe7eb2a79b47f9ffcab618cd611c2bc1ea6/FormalConjectures/ErdosProblems/700.lean}{Formal Conjectures, Erd\H{o}s Problem 700}.
Lean statements, revision \texttt{2c24ebe}.
\end{thebibliography}

\end{document}
